\chapter{Materiales y métodos}
\objetivo{Describir datos, reglas fijadas antes de leer el resultado, y el entorno, de forma que otra persona pueda repetir la corrida.}
\todojd{Lo que sigue es descriptivo y sale del preregistro y del certificado. Revísalo y reescríbelo con tus palabras. No cambies un umbral.}

\section{Cohortes}
\begin{table}[ht]
\caption{Matrices leídas por la corrida v3.7}
\label{tab:cohortes}
\begin{tabular}{p{3.2cm}p{2.6cm}p{7.6cm}}
\toprule
Matriz & Acceso & Papel en esta tesis \\
\midrule
Tejido TMT, CPTAC & PDC000116 & Tumor frente a normal publicado. Fuente: \texttt{config/datasets.yaml}. \cite{vasaikar2019} \\
Tejido LFQ & PXD034575 & Primario, metástasis hepática y mucosa. Fuente: \texttt{k1\_etiquetas}. \cite{tanaka2024} \\
Suero DIA & PXD055218 & Celda CP frente a CN. Fuente: \texttt{k4\_suero}. \cite{rai2025} \\
Suero de otro depósito & PXD064655 & No se calculó su contraste proteómico. H8 queda \HochoEstado. \\
\bottomrule
\end{tabular}
\end{table}
\fuente{\texttt{config/preregistro\_v3.yaml} y \texttt{config/datasets.yaml}.}

\section{Reglas de inferencia}
\begin{itemize}
\item Alfa 0,05. Benjamini-Hochberg. Margen \(\log_2(1{,}5)=0{,}585\). Mínimo de pares observados: 3. Mínimo por grupo en suero: 8. Fracción mínima observada: 0,70.
\item No se fusionan matrices. No se imputa. No se aplica un segundo \(\log_2\) a CPTAC.
\item Celda preregistrada del multiverso: PG.MaxLFQ, librería v3.3, paneles excluidos, mínimo \(n=3\), margen 0,585, alfa 0,05, sin normalización de tejido.
\item Clases del suero: diferente si TREAT pasa el corte en el sentido observado; equivalente si TOST cabe en el margen; indeterminada si queda fuera de las otras dos.
\end{itemize}
\todojd{Confirma en el código si los desplazamientos 720, 721, 840 y 841 se suman a la semilla base. El JSON de la corrida guarda la semilla \texttt{20260914} en \texttt{ml\_suero\_permutacion.json} y en \texttt{h7\_visibilidad\_en\_sangre.json}.}

\section{Modelos y priorización}
\begin{itemize}
\item M1 a M6, con los nombres de la enmienda: matriz humana, matriz principal, solo calidad, solo coagulación, principal residualizada por la calidad, y cinco componentes dentro del pliegue.
\item Validación escrita: 5 pliegues estratificados, 20 repeticiones, nulo de 1000 permutaciones.
\item Regla titular: si M2 no pasa la permutación, no hay separación detectable; si M2 y M5 pasan y el AUC de M5 es al menos 0,70, la señal va más allá de la calidad; en otro caso la separación se explica por la calidad.
\item Embudo de priorización, en este orden: \PrioriPasoUno, \PrioriPasoDos, \PrioriPasoTres, \PrioriPasoCuatro, \PrioriPasoCinco, \PrioriPasoSeis, \PrioriPasoSiete. Nadie pasa el paso 6.
\end{itemize}

\section{Lo que el método no hizo}
\begin{itemize}
\item Vías del suero, cláusula de vías de H1, unión de HC por XIC, deriva del PRTC, redes de coexpresión, catálogo de candidatas y neoantígenos. Están declarados como no ejecutados. No son un resultado.
\item H8: 0 tumorales epiteliales en un universo de \HochoUniverso. directLFQ 0.3.3 figura en el JSON. \cite{ammar2023}
\end{itemize}

\section{Entorno}
Python 3.14.3. numpy 2.4.4, pandas 3.0.2, scipy 1.17.1, scikit-learn 1.8.0. Semilla registrada en el sello: 20260914. Fuente: \texttt{results/v3\_7/certificado.json}.

\section{Uso de herramientas de inteligencia artificial}
\todojd{Ajusta este párrafo. La versión de Cursor y la fecha de SciSpace ya están en la declaración.}
Se utilizó Cursor (Anysphere) como asistente de programación para escribir y depurar código, y SciSpace como apoyo para buscar y hacer un primer filtro de literatura científica. El código, los resultados y las fuentes fueron revisados, ejecutados y validados por el autor, quien asume la responsabilidad total del contenido. \cite{cursor2026,scispace2026}

Cursor y SciSpace no se citan como fuente de un hecho científico.
Extensión orientativa: 8 a 12 páginas.
